% Reconstructed 2026-11-20 from complete authored text, final prism-interface deltas,
% and root/independent audits. New source identity; not historical byte identity.
\Needspace{10\baselineskip}
\subsection{A.156. Relative PL filling and the real weak comparison}

\begin{saresult}{A.156}{Proved under stated hypotheses}{Relative PL filling in every real base dimension}
Fix an o-minimal expansion of $\mathbb R$, a closed compact PL
manifold $B$, and $N=B\times I$. Let $P=[0,1]^m$, $m\ge1$,
and let $U_p:N\to N$ be a jointly continuous definable family
of homeomorphisms fixing $A_0=B\times\{0\}$ pointwise and
preserving $A_1=B\times\{1\}$ setwise. If its restriction to
$N\times\partial P$ is jointly PL, there is a jointly PL
family $W$ with the same spatial constraints and
\[
 W_p=U_p\quad(p\in\partial P).
 \tag{156.1}
\]
Consequently $P_{\rm PL}(B)\to P_{\rm def}(B)$ is a weak
equivalence for every such real $B$, including dimensions
three and four. This does not identify its comparison with
the whole ordinary TOP concordance space.
\end{saresult}

\begin{proof}
The first, spatial-clock part works over every o-minimal real
closed field. The later use of A.152 and the finite mesh in
the fragmentation part is over the real field.

\textbf{PL normalization of variation clocks.}\quad
For a fixed $D>3$, define, on $1\le d\le D$, $0\le v\le d$,
\[
 T_D(d,v)=
 \begin{cases}
 \min\{v,1+(v-d)/3\},&1\le d\le3,\\
 \min\{v/3,1+(v-d)/D\},&3\le d\le D.
 \end{cases}
 \tag{156.2}
\]
Both formulas equal $v/3$ at $d=3$. For fixed $d$ this is a
strictly increasing PL homeomorphism $[0,d]\to[0,1]$.
Its vertical slopes lie in $[1/D,1]$. It is globally
$1$-Lipschitz in the maximum norm: every affine gradient
has sum of absolute coefficients at most1, and its domain
is a convex trapezoid. In particular,
\[
 |v-v'|\le D|T_D(d,v)-T_D(d,v')|.
 \tag{156.3}
\]
At one coordinate stage of A.154 choose
$D>\max\{3,\sup_b A_b(1)\}$ and use the inverse of
\[
 \tau_b(t)=T_D(A_b(1),A_b(t))
 \tag{156.4}
\]
in place of the normalized raw clock. If $A_b(t)=3t$,
this clock is exactly the identity.

Here is the full mixed estimate. If two profiles $b,c$
differ by $\delta$ in a previously processed coordinate,
with old bound $K$, variation comparison gives
\[
 \sup_t|A_b(t)-A_c(t)|,\quad |A_b(1)-A_c(1)|
       \le C_0K\delta,\qquad C_0=2d_0N,
\]
where $d_0$ is the number of vector components and $N$ the
line monotonicity bound. Thus
$\sup_t|\tau_b(t)-\tau_c(t)|\le C_0K\delta$.
For $h=\tau_b^{-1}(u)$, $h'=\tau_c^{-1}(u)$,
\[
 |A_c(h)-A_c(h')|
 \le D|\tau_c(h)-\tau_c(h')|
 \le DC_0K\delta.
 \tag{156.5}
\]
The old-coordinate output bound is therefore
$(1+DC_0)K$, and the new-coordinate bound is $D$.
Including the evolving domain map among the vector
components, exactly as in A.154, proves uniform spatial
Lipschitz bounds for both final factors.

If a vector input is jointly PL on a parameter subpolyhedron,
its line variation is jointly PL there. Its exact finite
maximum formula uses linear constraints on the partition
points, sums, absolute values, and finite piecewise-affine
evaluation. Linear quantifier elimination gives a semilinear
graph for the maximum. Its continuity, proved in A.154,
therefore makes it a finite PL function. Hence the raw
clock, its endpoint height, (156.4), and the inverse total
clock are jointly PL on this parameter part. The inverse
is PL because it is the inverse of a PL homeomorphism of
the compact polyhedral product. Induction preserves this
property for every stage and both final factors.

The same argument respects coefficient fields. If boundary
data are PL over a subfield and $D$ belongs to that field,
all these boundary functions and inverses are PL over it.
On a unit-cube vector model one may choose $D$ to be a
standard integer larger than $3$ and $1+d_0N$.
This quantitative refinement is used later in A.160.

\textbf{Boundary factors and the fixed bottom.}\quad
Use the actual PL cubical map $b:N\to Q\subset C=[0,1]^n$
of A.154, compatible with both spatial ends. Thus
$U^Q=bUb^{-1}$ is jointly PL on $Q\times\partial P$.
Extend that restriction to a jointly PL vector map
$E_\partial:C\times\partial P\to C$: triangulate compatibly,
retain the prescribed vertex values, assign a fixed point
of $C$ at the other vertices, and extend affinely.
This is a vector extension, not a homeomorphic extension.

On $(Q\times P)\cup(C\times\partial P)$ prescribe $U^Q$
and $E_\partial$. The prescriptions agree. The explicit
full-subcomplex formula of A.154 gives a continuous definable
extension $E:C\times P\to C$. Apply (156.4) coordinatewise
to $(\mathrm{id},E)$. Restrict the resulting map to $Q$
and conjugate back. We obtain
\[
 U=s r^{-1},
 \tag{156.6}
\]
with uniform spatial Lipschitz bounds, with both factors
jointly PL on the parameter boundary, and preserving both
end marks.

The factors individually fix the bottom. It is a union
of coordinate faces on which $E=\mathrm{id}$. Inductively,
on a free coordinate line within such a face the evolving
vector is $(\mathrm{id},\mathrm{id})$, so its variation is
$2t$ and its raw clock is $3t$. Formula (156.4) is then
identity. A coordinate fixed by the face is an endpoint
and remains fixed. Thus every clock is pointwise identity
on every bottom face, and $r=s=\mathrm{id}$ on $A_0$.

The tagged graph family of these two factors has
Hausdorff-Lipschitz boundary, since its boundary maps are
jointly PL on a compact polyhedron. The exact relative
cube theorem A.152 gives $g:P\to P$, equal to identity on
the parameter boundary, with both pulled graph families
Hausdorff-Lipschitz. Their spatial bounds and the graph
matching estimate of A.154 give uniform forward parameter
Lipschitz bounds for
\[
 \bar r=r\circ g,\qquad \bar s=s\circ g.
 \tag{156.7}
\]
They are individually bottom-fixed and their parameter
boundary values remain exactly those of $r,s$.

\textbf{The extension-only relative fragmentation assertion.}\quad
We prove the assertion needed for each factor: a definable
concordance family $V$ which is uniformly forward
parameter-Lipschitz and jointly PL on $\partial P$ has
a jointly PL filling of that prescribed boundary.
No PL filling is assumed, and the intervening deformation
need not be fixed on the parameter boundary.

Double $N$ to $X=B\times S^1$, mark both halves and both
seams, and reflect $V$. Choose the Hudson functions and
finite target data of A.147 with $cL<1$. The fixed-point
and terminal support arguments use this inequality and
the spatial marks, not the identity value on parameter
faces. For the boundary track there is a necessary
normalization refinement. Put
\[
 \rho_{y,v}(q)=\min_{0\le j<R}\max\{f^j(y),q-2jv\},
 \quad T_v=1+2(R-1)v,\quad D=2R-1,
\]
and use
\[
 \Lambda(v,p)=\max\{p,D(p-1)+T_v\}.
 \tag{156.8}
\]
This is the jointly PL inverse, in the vertical coordinate,
of $(d,q)\mapsto(d,\min\{q,1+(q-d)/D\})$.
For fixed $v$ it takes $[0,1]$ homeomorphically to
$[0,T_v]$, and
\[
 \Lambda(0,p)=p,\quad\Lambda(1,p)=Dp,\quad
 \Lambda(v,0)=0,\quad\Lambda(v,1)=T_v.
\]
Use the fixed-point construction
\[
 t_i=\rho_{V_t(x),v}(\Lambda(v,p_i)),\qquad
 H_{p,v}(x)=V_t(x).
 \tag{156.9}
\]
Its contraction constant is unchanged. Its inverse is
\[
 H_{p,v}^{-1}(y)=
 V^{-1}_{(\rho_{y,v}(\Lambda(v,p_i)))_i}(y).
 \tag{156.10}
\]
Initial and terminal maps agree with those of A.147, since
the normalizations agree at $v=0,1$. In particular all
terminal support labels, strictification, and marked ball
coarsenings are unchanged.

The WHOLE parameter-boundary track is jointly PL.
If $p_i$ is0 or1, the corresponding coordinate in (156.10)
is that same endpoint. Thus the formula evaluates the
given jointly PL boundary inverse of $V$. Both $\rho$
and $\Lambda$ are jointly PL, so their composition with
that inverse is PL. The total boundary inverse map is a
PL homeomorphism, hence its inverse $H$ is PL as well.
This proves the assertion without treating the original
bilinear normalization $T_vp_i$ as PL.

\textbf{Prescribing the actual boundary in the PL endpoint.}\quad
Write $H=H_{-,1}$. The finite construction of A.147 gives
source and target prism decompositions over a triangulation
of $P$ compatible with $\partial P$, with cells
\[
 D_\sigma\times\sigma\longrightarrow C_\sigma\times\sigma.
 \tag{156.11}
\]
These are literal products over the entire closed simplex:
Mn\"ev's Proposition4, p.21, describes the target cells;
equation(11) and Proposition5, pp.22--23, give
\[
 H^{-1}(C\times\sigma)=H_{v_\sigma}^{-1}(C)\times\sigma.
\]
Thus each source spatial cell is fixed over the simplex.
All product boundaries are unions of proper lower-total-
dimensional prisms. Over a proper parameter face, spatial
cells may refine $D_\sigma$; the boundary conclusion is
unchanged. All spatial marks are subcomplexes.

On the parameter boundary, source spatial cells are
already PL: $D_\sigma=H_p^{-1}(C_\sigma)$ for any
$p\in\sigma$, and $H_p$ is PL there. Choose the fixed
target PL spatial complex to refine all target cells,
all these boundary source cells, and all marks.
Compatible triangulation and marked Hauptvermutung give
a definable homeomorphism $\phi$ preserving every one of
these PL cells setwise, with all other source cells PL
after taking inverse images.

The conjugate $\phi^{-1}H\phi$ has PL prism decompositions.
On a parameter-boundary prism, its source and target sets
are still the ORIGINAL $D_\sigma,C_\sigma$. Therefore
the ORIGINAL map $H|_{\partial P}$ is a jointly PL
cellular map with exactly the required correspondence.
Prescribe it on $X\times\partial P$ and prescribe identity
on $A_0\times P$. They agree on their intersection.
For the bottom correspondence, any target cell
$C\subset A_0$ has source $H_p^{-1}(C)=C$ because $H$
fixes the bottom pointwise, and $\phi$ preserves this
target cell. Thus the converted source is still $C$.

Rename each converted source cell $\phi^{-1}(D_\sigma)$
as $D_\sigma$. It is a PL ball: it is polyhedral and
definably homeomorphic, as a pair with boundary, to its
target ball; marked Hauptvermutung supplies a PL ball
chart. Products and incidence are retained by conjugation.

Order the unassigned prisms by total dimension. All of
\[
 \partial(D_\sigma\times\sigma)
 =(\partial D_\sigma\times\sigma)\cup
                      (D_\sigma\times\partial\sigma)
\]
has already been assigned a PL boundary homeomorphism
preserving the parameter. Choose compatible finite
subdivisions making it simplicial; refinements of earlier
triangulations change no assigned map or cell set.
Use fixed spatial PL charts to convex balls and cone
the TOTAL product boundary from $(d_0,b_\sigma)$ to
$(c_0,b_\sigma)$. The same base barycentre is used in
both balls. This is a PL homeomorphism and preserves
the parameter exactly: the base of corresponding points
is $(1-u)b_\sigma+up$. This is Mn\"ev's Lemma3,
proof Steps2--3, pp.24--25, not separate radial scaling
of a parameterized spatial map.
Finite induction gives a jointly PL family $G$, with
$G|_{\partial P}=H|_{\partial P}$, fixing the bottom
and preserving every mark. Restrict to the original half.

\textbf{Restoring the initial boundary by a PL collar.}\quad
Choose an affine inner cube $P_{\rm in}$ and a fixed PL
collar
\[
 c:\partial P\times I\longrightarrow
                    P\setminus\operatorname{int}P_{\rm in},
\]
whose inner edge is the corresponding affine copy of
$\partial P$ and whose outer edge is $\partial P$.
It can be made by triangulating the finite boundary
simplex prisms between homothetic cubes; no bilinear
radial formula is asserted to be PL.
Place a rescaled copy of $G$ on $P_{\rm in}$ and put
$H_{z,1-u}$ at $c(z,u)$. At the inner edge these agree
because $G_z=H_{z,1}$; at the outer edge the value is
$V_z$. The track was proved jointly PL, so this is
a jointly PL filling with exactly the initial boundary.
It retains the bottom and top constraints.

Apply the extension-only assertion separately to
$\bar r,\bar s$. It gives PL families $R,S$ with
boundary values $r,s$, respectively. Their quotient
$W=SR^{-1}$ is jointly PL, bottom-fixed, top-preserving,
and has boundary value $sr^{-1}=U$. This proves (156.1).

A based jointly PL $k$-cube nullhomotopic through based
definable concordances has a definable filling on a
$(k+1)$-cube with PL full boundary: the original cube
on one face and identity on the others. Formula (156.1)
gives a PL nullhomotopy. Hence comparison is injective
on every positive homotopy group. A.154 gives
surjectivity and A.113 gives the component bijection.
PL translation handles all base components.
Empty, disconnected and zero-dimensional bases are
treated as in the preceding marked constructions.
\end{proof}

